\section{Señal de programa final}
\label{sec:senal_programa}

Una vez procesada la mezcla, con el \textit{composite} activo, las transiciones aplicadas y
los \textit{overlays} superpuestos, \texttt{voctocore} emite la señal de programa final de forma
simultánea por dos puertos:

\begin{itemize}
  \item \textbf{Puerto 11000}: salida de la mezcla en bruto (vídeo sin comprimir en
        formato I420), utilizada principalmente por \textit{encoders} externos como FFmpeg para
        grabación local. La interfaz gráfica solo se conecta aquí cuando las
        previsualizaciones están desactivadas; en la configuración predeterminada
        utiliza el puerto~12000 (JPEG comprimido). Esta salida refleja siempre el
        estado real de la mezcla, independientemente del estado del \textit{stream blanker}.
  \item \textbf{Puerto 15000}: salida de \textit{streaming}, que pasa previamente por el
        \textit{stream blanker}. Cuando el estado es LIVE, esta salida es idéntica al puerto
        11000; cuando el estado es PAUSE o NOSTREAM, emite el vídeo alternativo
        correspondiente en lugar de la mezcla real.
\end{itemize}

El \textit{encoder} de salida (típicamente FFmpeg) se conecta al puerto 15000 y produce el
flujo H.264/AAC que se ingesta en la plataforma de distribución mediante los protocolos
\acs{RTMP} o \acs{SRT}, ampliamente soportados por las plataformas de \textit{streaming} en directo. La separación entre los puertos 11000 y
15000 garantiza que la grabación local refleje siempre la mezcla real del realizador,
mientras que la audiencia remota solo recibe la señal que el operador autoriza
explícitamente a través del \textit{stream blanker}.
